\begin{proof}[Proof of \cref{obj:thm:pilot-attainment}]
Fix the selector and pilot schedule in the statement, and let \(\mathcal P\) be the procedure in \cref{obj:def:pilot-estimator}. Use the staircase notation of \cref{obj:synth_4} and the contrast notation of \cref{obj:synth_6}; thus
\[
q=1-p,\qquad d=e^\varepsilon-1,\qquad c=(-1,1)^\top.
\]
Write
\[
N_n=\max\{n-m_n,0\}.
\]
Reciprocals of zero have the convention in \cref{obj:def:pilot-estimator,obj:synth_1}. For \(n\ge2\), the schedule assumptions give
\[
N_n=n-m_n>0,\qquad
N_n\longrightarrow\infty,\qquad
\frac{n}{N_n}\longrightarrow1.
\]
We use the local alternatives and scaled errors of \cref{obj:synth_2,obj:synth_3}, with baseline \(\theta\).

\begin{enumerate}
\item \emph{Bounds and centered conditional corrections.}
Define the reference weights, supported on the four singleton patterns, by
\[
\alpha^{\mathrm{RR}}_S=
\begin{cases}
(e^\varepsilon+3)^{-1},& |S|=1,\\
0,& |S|\ne1,
\end{cases}
\qquad S\in\mathcal S.
\]
Each record belongs to exactly one singleton pattern, so
\[
\sum_{S\in\mathcal S}\alpha^{\mathrm{RR}}_S b_S(j)
=\frac{e^\varepsilon+3}{e^\varepsilon+3}=1,
\qquad j\in\{0,1,2,3\}.
\]
Hence \(\alpha^{\mathrm{RR}}\in\mathcal A_\varepsilon\). For every interior parameter \(\vartheta\),
\[
1\le h_S(\vartheta)\le e^\varepsilon,
\qquad
0\le\alpha_S\le1
\quad(\alpha\in\mathcal A_\varepsilon).
\]
The first inequality follows because \(h_S(\vartheta)\) is an average of entries of \(b_S\), all of which lie in \([1,e^\varepsilon]\). The second follows from any channel-row normalization and \(b_S(j)\ge1\).

Define the constants
\[
\underline a_0=\frac{d^2q^2}{e^\varepsilon(e^\varepsilon+3)},
\qquad
\underline a_1=\frac{d^2p^2}{e^\varepsilon(e^\varepsilon+3)},
\qquad
\underline J=\frac{\underline a_0\underline a_1}
{\underline a_0+\underline a_1},
\]
and
\[
\overline F_0=\sum_{S\in\mathcal S}(g_S^\top v(0))^2,
\qquad
\overline t=\sqrt{\frac{\overline F_0}{\underline a_0}},
\qquad
\overline g=\sum_{S\in\mathcal S}(|g_{S,0}|+|g_{S,1}|),
\qquad
\overline\phi=\frac{\overline g(\overline t+1)}{\underline J}.
\]
Here \(\underline a_0,\underline a_1,\underline J>0\). The contributions of the singleton patterns \(\{0\}\) and \(\{3\}\) give, for every \(u\in\mathbb R\),
\[
F_\vartheta(\alpha^{\mathrm{RR}},u)
\ge \underline a_0u^2+\underline a_1(u+1)^2
=(\underline a_0+\underline a_1)
\left(u+\frac{\underline a_1}{\underline a_0+\underline a_1}\right)^2
+\underline J.
\]
Also, termwise use of \(\alpha_S\le1\) and \(h_S(\vartheta)\ge1\) gives
\[
F_\vartheta(\alpha,0)\le\overline F_0.
\]
The two saddle inequalities therefore imply
\[
\begin{aligned}
\underline J
&\le F_\vartheta(\alpha^{\mathrm{RR}},t_\vartheta)
\le F_\vartheta(\alpha_\vartheta,t_\vartheta)
=J_\vartheta\\
&\le F_\vartheta(\alpha_\vartheta,0)
\le\overline F_0,
\end{aligned}
\qquad
\underline a_0t_\vartheta^2\le J_\vartheta\le\overline F_0.
\]
Consequently,
\[
J_\vartheta=J^*(\vartheta,p,\varepsilon)\ge\underline J>0,
\qquad |t_\vartheta|\le\overline t.
\]

For \(\vartheta\in(0,1)^2\) and \(S\in\mathcal S\), define the correction at the selection parameter by
\[
\phi_\vartheta(S)
=\frac{g_S^\top v(t_\vartheta)}
{J_\vartheta h_S(\vartheta)}.
\]
At \(\vartheta=\widetilde\theta\), this is the correction in \cref{obj:synth_9}. Since
\[
|g_S^\top v(u)|
\le (|g_{S,0}|+|g_{S,1}|)(|u|+1)
\le\overline g(|u|+1),
\]
the preceding bounds yield
\[
|\phi_\vartheta(S)|\le\overline\phi
\qquad
(\vartheta\in(0,1)^2,\ S\in\mathcal S).
\]

We next establish the exact centering identities. Define
\[
\mathbf1_2=(1,1)^\top.
\]
Minimization in the profiling coordinate gives
\[
0=\left.\frac{d}{du}F_\vartheta(\alpha_\vartheta,u)
\right|_{u=t_\vartheta}
=2\mathbf1_2^\top I_\vartheta(\alpha_\vartheta)v(t_\vartheta).
\]
Thus \(I_\vartheta(\alpha_\vartheta)v(t_\vartheta)\) is a multiple of \(c\). Because
\[
c^\top v(t_\vartheta)=1,
\qquad
v(t_\vartheta)^\top I_\vartheta(\alpha_\vartheta)v(t_\vartheta)
=J_\vartheta,
\]
that multiple is \(J_\vartheta\), and hence
\[
I_\vartheta(\alpha_\vartheta)v(t_\vartheta)=J_\vartheta c.
\]
Channel normalization and the affine formula for \(h_S\) also give
\[
\sum_S\alpha_{\vartheta,S}g_S=0,
\qquad
h_S(\xi)=h_S(\vartheta)+g_S^\top(\xi-\vartheta).
\]

For \(\xi,\vartheta\in(0,1)^2\), define the main-release probabilities and centered correction by
\[
\omega_{\xi,\vartheta}(S)
=\alpha_{\vartheta,S}h_S(\xi),
\qquad
\zeta_{\xi,\vartheta}(S)
=\phi_\vartheta(S)-c^\top(\xi-\vartheta).
\]
These probabilities are nonnegative and sum to one. Expanding the affine denominator gives
\[
\begin{aligned}
\sum_S\omega_{\xi,\vartheta}(S)\phi_\vartheta(S)
&=\frac1{J_\vartheta}
\left\{
\sum_S\alpha_{\vartheta,S}g_S^\top v(t_\vartheta)
+(\xi-\vartheta)^\top
 I_\vartheta(\alpha_\vartheta)v(t_\vartheta)
\right\}\\
&=c^\top(\xi-\vartheta).
\end{aligned}
\]
At the diagonal,
\[
\sum_S\omega_{\vartheta,\vartheta}(S)\phi_\vartheta(S)^2
=\frac{F_\vartheta(\alpha_\vartheta,t_\vartheta)}{J_\vartheta^2}
=\frac1{J_\vartheta}
=V^*(\vartheta,p,\varepsilon).
\]
In particular,
\[
\sum_S\omega_{\xi,\vartheta}(S)\zeta_{\xi,\vartheta}(S)=0,
\qquad
|\zeta_{\xi,\vartheta}(S)|\le2\overline\phi.
\]
For the last inequality, the exact mean identity and the probability normalization imply
\[
|c^\top(\xi-\vartheta)|
=\left|\sum_S\omega_{\xi,\vartheta}(S)\phi_\vartheta(S)\right|
\le\overline\phi.
\]
% lean: CausalSmith.Stat.LdpAteEfficiencySurface.pilotSubsequenceRowModel

\item \emph{Continuity of the conditional variance.}
For \(\xi,\vartheta\in(0,1)^2\), define
\[
\sigma^2(\xi,\vartheta)
=\sum_S\omega_{\xi,\vartheta}(S)\zeta_{\xi,\vartheta}(S)^2
=\sum_S\alpha_{\vartheta,S}h_S(\xi)\phi_\vartheta(S)^2
-\{c^\top(\xi-\vartheta)\}^2.
\]
The diagonal second-moment identity gives
\[
\sigma^2(\vartheta,\vartheta)=V^*(\vartheta,p,\varepsilon).
\]
Subtracting this identity from the preceding formula yields
\[
\begin{aligned}
\sigma^2(\xi,\vartheta)-V^*(\vartheta,p,\varepsilon)
={}&\sum_S\alpha_{\vartheta,S}
\{h_S(\xi)-h_S(\vartheta)\}\phi_\vartheta(S)^2\\
&-\{c^\top(\xi-\vartheta)\}^2.
\end{aligned}
\]
Therefore,
\[
\left|\sigma^2(\xi,\vartheta)-V^*(\vartheta,p,\varepsilon)\right|
\le
\overline\phi^2\sum_S|h_S(\xi)-h_S(\vartheta)|
+\{c^\top(\xi-\vartheta)\}^2.
\]

To establish continuity of the optimized value, define
\[
K=[-\overline t,\overline t],
\qquad
\mathcal E_\xi(u)=\max_{\alpha\in\mathcal A_\varepsilon}
F_\xi(\alpha,u),
\quad \xi\in(0,1)^2,\ u\in\mathbb R.
\]
The saddle inequalities give
\[
\mathcal E_\xi(t_\xi)=J_\xi,
\qquad
\mathcal E_\xi(u)\ge
F_\xi(\alpha_\xi,u)\ge J_\xi.
\]
Since \(t_\xi\in K\),
\[
J^*(\xi,p,\varepsilon)=\min_{u\in K}\mathcal E_\xi(u).
\]
The feasible set is nonempty, closed, and contained in \([0,1]^{14}\), hence compact. The objective is jointly continuous on interior parameters, feasible weights, and \(K\). The elementary bounds for maxima and minima give
\[
\left|J^*(\xi,p,\varepsilon)-J^*(\theta,p,\varepsilon)\right|
\le
\sup_{\substack{\alpha\in\mathcal A_\varepsilon\\u\in K}}
|F_\xi(\alpha,u)-F_\theta(\alpha,u)|
\longrightarrow0
\quad(\xi\to\theta).
\]
The convergence follows from joint continuity on a compact parameter neighborhood times \(\mathcal A_\varepsilon\times K\). Positivity of \(J^*\) then gives
\[
V^*(\xi,p,\varepsilon)\longrightarrow V^*(\theta,p,\varepsilon).
\]
Combining this with the displayed variance-error bound proves
\[
\sigma^2(\xi,\vartheta)\longrightarrow V^*(\theta,p,\varepsilon)
\qquad
\text{as }(\xi,\vartheta)\to(\theta,\theta)
\text{ through interior pairs}.
\]
% lean: CausalSmith.Stat.LdpAteEfficiencySurface.continuousAt_adaptiveScoreVariance_diag

\item \emph{Gaussian convergence of deterministic conditional rows.}
Fix \(h\in\mathbb R^2\). To define valid row probabilities at every index, set
\[
\xi_{n,h}=
\begin{cases}
\theta_{n,h},&\theta_{n,h}\in(0,1)^2,\\
\theta,&\theta_{n,h}\notin(0,1)^2.
\end{cases}
\]
By \cref{obj:synth_2}, \(\theta_{n,h}\to\theta\); thus it is eventually interior, and
\[
\xi_{n,h}\to\theta,
\qquad
\xi_{n,h}=\theta_{n,h}\quad\text{eventually}.
\]

For \(n\ge2\) and \(\vartheta\in(0,1)^2\), let
\(\mathbb P^{\mathrm{row}}_{n,h,\vartheta}\) be the product law of \(N_n\) independent pattern releases with
\[
\mathbb P^{\mathrm{row}}_{n,h,\vartheta}(S_j=S)
=\omega_{\xi_{n,h},\vartheta}(S),
\qquad 1\le j\le N_n,\quad S\in\mathcal S.
\]
Define its total-sample-scaled distribution function by
\[
\mathcal C_{n,h}(x;\vartheta)
=
\mathbb P^{\mathrm{row}}_{n,h,\vartheta}
\left\{
\frac{\sqrt n}{N_n}
\sum_{j=1}^{N_n}\zeta_{\xi_{n,h},\vartheta}(S_j)\le x
\right\},
\qquad x\in\mathbb R.
\]
Also define the target distribution function
\[
\mathcal G_\theta(x)
=\Phi\left(\frac{x}{\sqrt{V^*(\theta,p,\varepsilon)}}\right).
\]

Take any strictly increasing sequence of integers \(n_k\ge2\) and any interior sequence \(\vartheta_k\to\theta\). Under the corresponding row law, put
\[
Z_{k,j}=\zeta_{\xi_{n_k,h},\vartheta_k}(S_j),
\qquad 1\le j\le N_{n_k},
\]
and
\[
\mathcal D_k(y)
=
\mathbb P^{\mathrm{row}}_{n_k,h,\vartheta_k}
\left\{
N_{n_k}^{-1/2}\sum_{j=1}^{N_{n_k}}Z_{k,j}\le y
\right\},
\qquad y\in\mathbb R.
\]
The rows are independent, centered, and bounded by \(2\overline\phi\), and their single-release variances satisfy
\[
\mathbb E^{\mathrm{row}}[Z_{k,1}^2]
=\sigma^2(\xi_{n_k,h},\vartheta_k)
\longrightarrow V^*(\theta,p,\varepsilon)>0.
\]
The increments normalized by
\(\sqrt{N_{n_k}V^*(\theta,p,\varepsilon)}\) therefore satisfy
\[
\sum_{j=1}^{N_{n_k}}
\mathbb E^{\mathrm{row}}
\left[
\left(\frac{Z_{k,j}}
{\sqrt{N_{n_k}V^*(\theta,p,\varepsilon)}}\right)^2
\right]
=
\frac{\sigma^2(\xi_{n_k,h},\vartheta_k)}
{V^*(\theta,p,\varepsilon)}
\longrightarrow1,
\]
and
\[
\sum_{j=1}^{N_{n_k}}
\mathbb E^{\mathrm{row}}
\left[
\left|\frac{Z_{k,j}}
{\sqrt{N_{n_k}V^*(\theta,p,\varepsilon)}}\right|^4
\right]
\le
\frac{(2\overline\phi)^4}
{N_{n_k}\{V^*(\theta,p,\varepsilon)\}^2}
\longrightarrow0.
\]
The triangular-array Lyapunov central limit theorem consequently gives
\[
\mathcal D_k(y)\longrightarrow\mathcal G_\theta(y)
\qquad(y\in\mathbb R).
\]

Define
\[
\lambda_k=\sqrt{\frac{n_k}{N_{n_k}}},
\qquad
\chi_k=\frac{x}{\lambda_k}.
\]
Then \(\lambda_k\to1\), \(\chi_k\to x\), and
\[
\mathcal C_{n_k,h}(x;\vartheta_k)=\mathcal D_k(\chi_k).
\]
For every \(a>0\), eventually
\[
\mathcal D_k(x-a)\le\mathcal D_k(\chi_k)\le\mathcal D_k(x+a).
\]
Pointwise convergence and continuity of \(\mathcal G_\theta\), followed by \(a\downarrow0\), prove
\[
\mathcal C_{n_k,h}(x;\vartheta_k)\longrightarrow\mathcal G_\theta(x).
\]
% lean: CausalSmith.Stat.LdpAteEfficiencySurface.adaptivePilotRowScaledCDF_subsequence_tendsto_gaussian

\item \emph{Averaging over the private pilot.}
For \(n\ge2\) and a pilot release vector
\(\mathbf r=(r_1,\ldots,r_{m_n})\in\{0,1,2,3\}^{m_n}\), define
\[
\vartheta_n(\mathbf r)
=\widetilde\theta
\quad\text{computed from }R_i=r_i
\text{ by }\cref{obj:def:pilot-estimator}.
\]
Clipping ensures \(\vartheta_n(\mathbf r)\in(0,1)^2\) for every \(\mathbf r\). Define the pilot weights
\[
w_{n,h}(\mathbf r)
=
\prod_{i=1}^{m_n}
\frac{1+d\,\pi_{\xi_{n,h},r_i}}{e^\varepsilon+3}.
\]
By \cref{obj:synth_10}, these weights are nonnegative and sum to one. Factorization of the successive release kernels gives, whenever \(\theta_{n,h}\) is interior,
\[
\Pr_{\theta_{n,h}}\{U_{n,h}^{\mathcal P}\le x\}
=
\sum_{\mathbf r}w_{n,h}(\mathbf r)
\mathcal C_{n,h}(x;\vartheta_n(\mathbf r)).
\]
Indeed, the exact mean identity gives, for each fixed pilot,
\[
U_{n,h}^{\mathcal P}
=
\frac{\sqrt n}{N_n}
\sum_{j=1}^{N_n}
\zeta_{\theta_{n,h},\vartheta_n(\mathbf r)}(S_j).
\]

We record the pilot concentration bound used in this averaging. Define
\[
\ell_0=\frac{qd}{e^\varepsilon+3},
\qquad
\ell_1=\frac{pd}{e^\varepsilon+3},
\qquad
\mathcal B_n(\varrho)
=
\frac1{m_n(\varrho\ell_0)^2}
+\frac1{m_n(\varrho\ell_1)^2},
\quad \varrho>0,\ n\ge2.
\]
For an interior parameter \(\xi\), the means of the pilot indicators for categories \(1\) and \(3\) are
\[
\frac{1+dq\xi_0}{e^\varepsilon+3},
\qquad
\frac{1+dp\xi_1}{e^\varepsilon+3}.
\]
Each indicator has variance at most one, so independence and Chebyshev's inequality give, for \(k\in\{1,3\}\) and \(b>0\),
\[
\Pr_\xi\{|\widehat u_k-\mathbb E_\xi\widehat u_k|\ge b\}
\le\frac1{m_nb^2}.
\]
Frequency inversion scales the two coordinate errors by
\(\ell_0^{-1}\) and \(\ell_1^{-1}\). If
\[
\delta_n+\varrho\le
\min\{\xi_0,1-\xi_0,\xi_1,1-\xi_1\},
\]
then an inverted coordinate within \(\varrho\) of \(\xi\) lies inside the clipping interval. Taking the contrapositive coordinatewise and applying the union bound yields
\[
\Pr_\xi\{\|\widetilde\theta-\xi\|_\infty\ge\varrho\}
\le\mathcal B_n(\varrho),
\qquad
\|\widetilde\theta-\xi\|_\infty
=\max_{k\in\{0,1\}}|\widetilde\theta_k-\xi_k|.
\]

Define, for \(n\ge2\),
\[
\beta_n=m_n^{-1/4},
\qquad
\varrho_{n,h}
=
2\max\left\{
\beta_n,\sqrt{\|\theta_{n,h}-\theta\|_\infty}
\right\},
\]
and
\[
\Delta_\theta
=\min\{\theta_0,1-\theta_0,\theta_1,1-\theta_1\}>0.
\]
Pilot divergence and local convergence imply
\[
\beta_n\to0,\qquad \varrho_{n,h}\to0.
\]
Eventually,
\[
\beta_n<\Delta_\theta/4,\qquad
\delta_n<\Delta_\theta/4,\qquad
\|\theta_{n,h}-\theta\|_\infty<\Delta_\theta/4,
\]
so the concentration bound applies at \(\xi=\theta_{n,h}\) with radius \(\beta_n\). Its right-hand side becomes
\[
\mathcal B_n(\beta_n)
=
\frac1{\sqrt{m_n}}
\left(\frac1{\ell_0^2}+\frac1{\ell_1^2}\right)
\longrightarrow0.
\]
Moreover, eventually
\[
\|\theta_{n,h}-\theta\|_\infty
\le\sqrt{\|\theta_{n,h}-\theta\|_\infty}.
\]
Hence
\[
\|\vartheta_n(\mathbf r)-\theta_{n,h}\|_\infty<\beta_n
\quad\Longrightarrow\quad
\|\vartheta_n(\mathbf r)-\theta\|_\infty\le\varrho_{n,h}.
\]

The deterministic-row convergence is uniform over these shrinking pilot neighborhoods. Explicitly, for every \(a>0\), eventually
\[
\|\vartheta_n(\mathbf r)-\theta\|_\infty\le\varrho_{n,h}
\quad\Longrightarrow\quad
|\mathcal C_{n,h}(x;\vartheta_n(\mathbf r))
-\mathcal G_\theta(x)|<a.
\]
If this implication failed, one could choose strictly increasing \(n_k\) and pilot vectors \(\mathbf r_k\) such that
\[
\|\vartheta_{n_k}(\mathbf r_k)-\theta\|_\infty
\le\varrho_{n_k,h}\longrightarrow0,
\qquad
|\mathcal C_{n_k,h}(x;\vartheta_{n_k}(\mathbf r_k))
-\mathcal G_\theta(x)|\ge a.
\]
This contradicts the deterministic-row limit.

Both distribution functions lie in \([0,1]\). Splitting the pilot average into the displayed neighborhood and its complement therefore gives, eventually,
\[
\left|
\sum_{\mathbf r}w_{n,h}(\mathbf r)
\mathcal C_{n,h}(x;\vartheta_n(\mathbf r))
-\mathcal G_\theta(x)
\right|
\le a+\mathcal B_n(\beta_n).
\]
Letting \(n\to\infty\), then \(a\downarrow0\), proves
\[
\Pr_{\theta_{n,h}}\{U_{n,h}^{\mathcal P}\le x\}
\longrightarrow\mathcal G_\theta(x)
\qquad(h\in\mathbb R^2,\ x\in\mathbb R).
\]
% lean: CausalSmith.Stat.LdpAteEfficiencySurface.pilot_scaledError_cdf_tendsto_of_all_conditionalRows

\item \emph{The weak local limit.}
For an interior parameter \(\xi\), the transcript law is a probability law: it is a mixture of conditional transcript probability laws, with nonnegative input-path weights satisfying
\[
\sum_{(j_1,\ldots,j_n)\in\{0,1,2,3\}^n}
\prod_{i=1}^n\pi_{\xi,j_i}
=\left(\sum_{j=0}^3\pi_{\xi,j}\right)^n=1.
\]
This also holds at \(n=0\), with the empty-product convention.

For each fixed \(h\), define a Borel probability law on \(\mathbb R\) by
\[
\nu_{n,h}=
\begin{cases}
\mathcal L_{\theta_{n,h}}(U_{n,h}^{\mathcal P}),
&\theta_{n,h}\in(0,1)^2,\\
N(0,V^*(\theta,p,\varepsilon)),
&\theta_{n,h}\notin(0,1)^2.
\end{cases}
\]
The local parameter is eventually interior, so the preceding distribution-function limit applies to \(\nu_{n,h}\). Convergence at every point of the continuous Gaussian distribution function implies weak convergence. Equivalently, for every bounded continuous \(f:\mathbb R\to\mathbb R\),
\[
\int f(u)\,d\nu_{n,h}(u)
\longrightarrow
\int f(u)\,dN(0,V^*(\theta,p,\varepsilon))(u).
\]
Eventually the left-hand side is
\(\mathbb E_{\theta_{n,h}}[f(U_{n,h}^{\mathcal P})]\), proving the asserted common weak local limit. Taking \(h=0\) gives the stated convergence at \(\theta\).
% lean: CausalSmith.Stat.LdpAteEfficiencySurface.weakLocalLimit_of_scaledError_cdf_tendsto

\item \emph{Consistency of the empirical variance.}
The empirical variance is the quantity in \cref{obj:synth_1}. Fix \(\delta>0\). Continuity of \(\sigma^2\) at \((\theta,\theta)\) supplies \(R_{\mathrm{var}}>0\) such that
\[
\|\vartheta-\theta\|_\infty<R_{\mathrm{var}},
\quad \vartheta\in(0,1)^2
\quad\Longrightarrow\quad
|\sigma^2(\theta,\vartheta)-V^*(\theta,p,\varepsilon)|<\delta/2.
\]
Set
\[
\varrho_{\mathrm{var}}
=\min\{R_{\mathrm{var}}/2,\Delta_\theta/4\}>0.
\]
Eventually \(\delta_n<\Delta_\theta/2\), so
\[
\delta_n+\varrho_{\mathrm{var}}\le\Delta_\theta,
\qquad
\Pr_\theta\{\|\widetilde\theta-\theta\|_\infty
\ge\varrho_{\mathrm{var}}\}
\le\mathcal B_n(\varrho_{\mathrm{var}})\longrightarrow0.
\]

Condition on a pilot vector, and write its estimate as \(\vartheta\). Define the centered main-sample moments
\[
\overline\zeta_n
=\frac1{N_n}\sum_{j=1}^{N_n}\zeta_{\theta,\vartheta}(S_j),
\qquad
\overline{\zeta^2}_n
=\frac1{N_n}\sum_{j=1}^{N_n}\zeta_{\theta,\vartheta}(S_j)^2.
\]
Translation invariance of empirical variance gives
\[
\widehat V^*=\overline{\zeta^2}_n-\overline\zeta_n^{\,2}.
\]
Conditional independence, centering, and boundedness give
\[
\mathbb E_\theta[\overline\zeta_n^{\,2}\mid H_{m_n}]
=\frac{\sigma^2(\theta,\vartheta)}{N_n}
\le\frac{(2\overline\phi)^2}{N_n},
\]
and
\[
\operatorname{Var}_\theta(\overline{\zeta^2}_n\mid H_{m_n})
\le\frac{(2\overline\phi)^4}{N_n}.
\]
On a pilot with
\(|\sigma^2(\theta,\vartheta)-V^*(\theta,p,\varepsilon)|<\delta/2\),
the event \(|\widehat V^*-V^*(\theta,p,\varepsilon)|>\delta\) implies
\[
|\overline{\zeta^2}_n-\sigma^2(\theta,\vartheta)|\ge\delta/4
\quad\text{or}\quad
\overline\zeta_n^{\,2}\ge\delta/4.
\]
Chebyshev's and Markov's inequalities consequently bound its conditional probability by
\[
\frac{(2\overline\phi)^4}{N_n(\delta/4)^2}
+\frac{(2\overline\phi)^2}{N_n(\delta/4)}.
\]
On the complementary pilots the conditional probability is at most one. Averaging gives, eventually,
\[
\begin{aligned}
\Pr_\theta\{|\widehat V^*-V^*(\theta,p,\varepsilon)|>\delta\}
\le{}&
\mathcal B_n(\varrho_{\mathrm{var}})
+\frac{(2\overline\phi)^4}{N_n(\delta/4)^2}\\
&+\frac{(2\overline\phi)^2}{N_n(\delta/4)}
\longrightarrow0.
\end{aligned}
\]
% lean: CausalSmith.Stat.LdpAteEfficiencySurface.pilot_VhatStar_consistent

\item \emph{Uniform squared tails and regularity.}
Fix \(H>0\), and use the admissible perturbation set of \cref{obj:synth_6}. For sufficiently large \(n\),
\[
\frac n{N_n}\le2.
\]
For every admissible \(h\), condition on the pilot and denote the \(N_n\) centered corrections by
\[
Z_j=\zeta_{\theta_{n,h},\widetilde\theta}(S_j),
\qquad 1\le j\le N_n.
\]
They are conditionally independent and identically distributed, with mean zero and
\[
|Z_j|\le2\overline\phi,
\qquad
U_{n,h}^{\mathcal P}
=\frac{\sqrt n}{N_n}\sum_{j=1}^{N_n}Z_j.
\]
Define their conditional moments by
\[
\mathfrak m_2=\mathbb E_{\theta_{n,h}}[Z_1^2\mid H_{m_n}],
\qquad
\mathfrak m_4=\mathbb E_{\theta_{n,h}}[Z_1^4\mid H_{m_n}].
\]
Then
\[
0\le\mathfrak m_2\le(2\overline\phi)^2,
\qquad
0\le\mathfrak m_4\le(2\overline\phi)^4.
\]
In the fourth-power expansion, terms having an index that occurs exactly once have zero conditional expectation. The remaining terms give
\[
\mathbb E_{\theta_{n,h}}
\left[\left(\sum_{j=1}^{N_n}Z_j\right)^4\middle|H_{m_n}\right]
=
N_n\mathfrak m_4+3N_n(N_n-1)\mathfrak m_2^2.
\]
Since \(N_n\ge1\),
\[
\mathbb E_{\theta_{n,h}}
\left[\left(N_n^{-1/2}\sum_{j=1}^{N_n}Z_j\right)^4
\middle|H_{m_n}\right]
\le4(2\overline\phi)^4.
\]
Thus, uniformly over the pilot and admissible \(h\),
\[
\mathbb E_{\theta_{n,h}}
[(U_{n,h}^{\mathcal P})^4\mid H_{m_n}]
\le
4\left(\frac n{N_n}\right)^2(2\overline\phi)^4
\le16(2\overline\phi)^4.
\]
Averaging preserves this bound. For every integer \(M>0\), the pointwise inequality
\[
u^2\mathbf1\{|u|>M\}\le\frac{u^4}{M^2},
\qquad u\in\mathbb R,
\]
therefore yields
\[
0\le
\limsup_{n\to\infty}
\sup_{h\in\mathcal H_{n,H}(\theta)}
\mathbb E_{\theta_{n,h}}
\left[(U_{n,h}^{\mathcal P})^2
\mathbf1\{|U_{n,h}^{\mathcal P}|>M\}\right]
\le\frac{16(2\overline\phi)^4}{M^2}.
\]
Letting \(M\to\infty\) proves the required uniform squared-tail convergence.

The Gaussian law
\[
L_{\mathcal P}=N(0,V^*(\theta,p,\varepsilon))
\]
is a probability law with
\[
\int u^2\,dL_{\mathcal P}(u)<\infty.
\]
Together with the common weak local limit and the squared-tail bound, this verifies the regularity conditions in \cref{obj:synth_6}.
% lean: CausalSmith.Stat.LdpAteEfficiencySurface.regular_of_weakLocalLimit_and_uniformTails

\item \emph{Standardizing the limiting distribution function.}
The reference-design argument gives
\[
J^*(\theta,p,\varepsilon)>0,
\qquad
V^*(\theta,p,\varepsilon)>0,
\qquad
\sqrt{V^*(\theta,p,\varepsilon)}>0.
\]
Define, for each sample size,
\[
\mathsf T_n
=\frac{U_{n,0}^{\mathcal P}}
{\sqrt{V^*(\theta,p,\varepsilon)}}.
\]
Since \(\theta_{n,0}=\theta\), positivity of the denominator gives, for every \(x\in\mathbb R\),
\[
\{\mathsf T_n\le x\}
=
\left\{U_{n,0}^{\mathcal P}
\le x\sqrt{V^*(\theta,p,\varepsilon)}\right\}.
\]
The Gaussian scaling identity is
\[
N(0,V^*(\theta,p,\varepsilon))
\left((-\infty,x\sqrt{V^*(\theta,p,\varepsilon)}]\right)
=\Phi(x).
\]
Applying the local distribution-function limit at this threshold proves
\[
\Pr_\theta\{\mathsf T_n\le x\}\longrightarrow\Phi(x)
\qquad(x\in\mathbb R).
\]
% lean: CausalSmith.Stat.LdpAteEfficiencySurface.pilot_attainment

\item \emph{Transfer to the Wald pivot.}
First, consistency also controls non-strict deviation events. For every \(b>0\),
\[
\{|\widehat V^*-V^*(\theta,p,\varepsilon)|\ge b\}
\subseteq
\{|\widehat V^*-V^*(\theta,p,\varepsilon)|>b/2\},
\]
and hence
\[
\Pr_\theta\{|\widehat V^*-V^*(\theta,p,\varepsilon)|\ge b\}
\longrightarrow0.
\]
The normal distribution function is continuous and strictly increasing, with limits zero and one at the two ends of the real line. The critical-value definition thus gives
\[
\Phi(z_{1-\gamma/2})=1-\gamma/2.
\]
Normal symmetry gives \(\Phi(0)=1/2\). Since \(1-\gamma/2>1/2\),
\[
z_{1-\gamma/2}>0.
\]

Define the estimation error and fallback pivot, for \(n\ge1\), by
\[
\Delta_{\tau,n}=\widehat\tau^*-\tau,
\qquad
\mathsf W_n=
\begin{cases}
0,&\widehat V^*=0,\ \Delta_{\tau,n}=0,\\
z_{1-\gamma/2}+1,&\widehat V^*=0,\ \Delta_{\tau,n}\ne0,\\
\sqrt n\,\Delta_{\tau,n}/\sqrt{\widehat V^*},
&\widehat V^*>0.
\end{cases}
\]
The empirical variance is nonnegative on every transcript, since it is a nonnegative multiple of a sum of squares. When it is positive, ordinary square-root algebra gives
\[
|\mathsf W_n|\le z_{1-\gamma/2}
\quad\Longleftrightarrow\quad
|\Delta_{\tau,n}|
\le z_{1-\gamma/2}\sqrt{\widehat V^*/n}.
\]
When it is zero, both events hold precisely when \(\Delta_{\tau,n}=0\), by the two fallback branches. Thus the equivalence holds on every transcript.

We show that \(\mathsf W_n-\mathsf T_n\) converges to zero in probability. The standardized distribution-function limit first gives tightness. Explicitly, for any \(\eta>0\), define
\[
\eta_{\mathrm{tail}}=\min\{\eta/8,1/8\}.
\]
Choose \(R_{\mathrm{tight}}>0\) so that the standard normal probabilities of
\((-\infty,-R_{\mathrm{tight}}]\) and
\([R_{\mathrm{tight}},\infty)\) are each at most
\(\eta_{\mathrm{tail}}\). Distribution-function convergence gives weak convergence of \(\mathsf T_n\), and the two endpoints have zero normal mass. Consequently,
\[
\Pr_\theta\{\mathsf T_n\notin[-R_{\mathrm{tight}},R_{\mathrm{tight}}]\}
\longrightarrow
1-\Phi(R_{\mathrm{tight}})+\Phi(-R_{\mathrm{tight}})
\le2\eta_{\mathrm{tail}}<\eta.
\]
Taking
\[
M_{\mathrm{tight}}=R_{\mathrm{tight}}+1
\]
therefore gives, eventually,
\[
\Pr_\theta\{M_{\mathrm{tight}}\le|\mathsf T_n|\}<\eta.
\]

Now fix \(\delta>0\) and \(\eta>0\), and choose such a positive
\(M_{\mathrm{tight}}\) with tail probability eventually below \(\eta/2\).
Define
\[
\mathfrak g_{\mathrm{scale}}(u)
=\frac{\sqrt{V^*(\theta,p,\varepsilon)}}{\sqrt u},
\qquad u>0.
\]
This function is continuous at \(V^*(\theta,p,\varepsilon)\), where it equals one. Choose \(\varrho_{\mathrm{scale},0}>0\) so that
\[
u>0,\quad |u-V^*(\theta,p,\varepsilon)|<\varrho_{\mathrm{scale},0}
\quad\Longrightarrow\quad
|\mathfrak g_{\mathrm{scale}}(u)-1|
<\delta/M_{\mathrm{tight}},
\]
and set
\[
\varrho_{\mathrm{scale}}
=
\min\{\varrho_{\mathrm{scale},0}/2,
V^*(\theta,p,\varepsilon)/2\}>0.
\]
On the event
\[
|\widehat V^*-V^*(\theta,p,\varepsilon)|
<\varrho_{\mathrm{scale}},
\qquad
|\mathsf T_n|<M_{\mathrm{tight}},
\]
the estimated variance is positive and
\[
\mathsf W_n
=\mathsf T_n\mathfrak g_{\mathrm{scale}}(\widehat V^*),
\qquad
|\mathsf W_n-\mathsf T_n|
\le M_{\mathrm{tight}}
|\mathfrak g_{\mathrm{scale}}(\widehat V^*)-1|
<\delta.
\]
It follows that
\[
\begin{aligned}
\Pr_\theta\{|\mathsf W_n-\mathsf T_n|\ge\delta\}
\le{}&
\Pr_\theta\{|\widehat V^*-V^*(\theta,p,\varepsilon)|
\ge\varrho_{\mathrm{scale}}\}\\
&+\Pr_\theta\{|\mathsf T_n|\ge M_{\mathrm{tight}}\}
<\eta
\end{aligned}
\]
eventually. Thus this deviation probability tends to zero.

For every \(x\in\mathbb R\) and \(\delta>0\),
\[
\begin{aligned}
\Pr_\theta\{\mathsf T_n\le x-\delta\}
-\Pr_\theta\{|\mathsf W_n-\mathsf T_n|\ge\delta\}
&\le\Pr_\theta\{\mathsf W_n\le x\}\\
&\le\Pr_\theta\{\mathsf T_n\le x+\delta\}
+\Pr_\theta\{|\mathsf W_n-\mathsf T_n|\ge\delta\}.
\end{aligned}
\]
Taking limits and then letting \(\delta\downarrow0\) proves convergence of the pivot distribution functions to \(\Phi\). The closed interval
\([-z_{1-\gamma/2},z_{1-\gamma/2}]\) has zero normal boundary mass, so
\[
\Pr_\theta\{|\mathsf W_n|\le z_{1-\gamma/2}\}
\longrightarrow
N(0,1)([-z_{1-\gamma/2},z_{1-\gamma/2}]).
\]
The event equivalence transfers this limit to the stated Wald event.
% lean: CausalSmith.Stat.LdpAteEfficiencySurface.pilot_wald_coverage_of_standardized_cdf_and_VhatStar_consistent

\item \emph{Privacy at every sample size and history.}
The randomized-response probabilities in \cref{obj:synth_10} are normalized, and for every pair of private records and every reported category,
\[
Q_{\mathrm{RR}}(k\mid x_j)
\le e^\varepsilon Q_{\mathrm{RR}}(k\mid x_{j'}).
\]
Summing over an output event proves the pilot privacy inequality.

At any main stage, including at an unattainable history, the history-based parameter estimate clips both coordinates to
\[
[\delta_n,1-\delta_n]\subset(0,1),
\qquad
0<\delta_n=(m_n+2)^{-1}\le1/2.
\]
The selected weights are therefore feasible. For every pattern \(S\),
\[
Q_{\alpha_\vartheta}(S\mid x_j)
=\alpha_{\vartheta,S}b_S(j),
\qquad
\sum_SQ_{\alpha_\vartheta}(S\mid x_j)=1,
\]
and
\[
\alpha_{\vartheta,S}b_S(j)
\le e^\varepsilon\alpha_{\vartheta,S}b_S(j').
\]
This inequality also holds for zero weights. Summing over patterns proves privacy for every measurable main-release event. At each fixed history, the same selected channel is used on the two compared private records. The two stage cases therefore verify the Markov property and the all-history inequality in \cref{obj:ass:sequential-ldp}. At \(n=0\), the stage conditions are vacuous.
% lean: CausalSmith.Stat.LdpAteEfficiencySurface.pilotProtocol_private

\item \emph{Exact pilot-adapted unbiasedness.}
Fix \(n\ge2\). Factorization gives, for a pilot vector \(\mathbf r\) and a main pattern vector
\(\mathbf s=(s_1,\ldots,s_{N_n})\in\mathcal S^{N_n}\),
\[
\Pr_\theta\{H_{m_n}=\mathbf r,\ (S_{m_n+1},\ldots,S_n)=\mathbf s\}
=
w_{n,0}(\mathbf r)
\prod_{j=1}^{N_n}
\omega_{\theta,\vartheta_n(\mathbf r)}(s_j).
\]
Transcripts with a release type incompatible with its stage have zero mass.

For each fixed pilot, product normalization and the exact correction mean give
\[
\begin{aligned}
&\sum_{\mathbf s\in\mathcal S^{N_n}}
\left(\prod_{j=1}^{N_n}
\omega_{\theta,\vartheta_n(\mathbf r)}(s_j)\right)
\left\{
c^\top\vartheta_n(\mathbf r)
+\frac1{N_n}\sum_{j=1}^{N_n}
\phi_{\vartheta_n(\mathbf r)}(s_j)
-c^\top\theta
\right\}\\
&\hspace{2cm}
=c^\top\vartheta_n(\mathbf r)
+\frac{N_n}{N_n}
c^\top(\theta-\vartheta_n(\mathbf r))
-c^\top\theta
=0.
\end{aligned}
\]
Multiplication by \(w_{n,0}(\mathbf r)\) shows that the integral of
\(\widehat\tau^*-\tau\) over every pilot-prefix cell is zero. A prefix incompatible with the pilot release type has zero mass, so its cell integral is also zero.

A function \(f\) satisfying the statement is constant on each such cell. Summing the zero cell integrals, each multiplied by its corresponding value of \(f\), gives
\[
\mathbb E_\theta[(\widehat\tau^*-\tau)f]=0.
\]
The transcript space is finite and the estimator and \(f\) are real-valued, so their product is integrable under the transcript probability law.
% lean: CausalSmith.Stat.LdpAteEfficiencySurface.pilot_conditional_unbiasedness

\item \emph{Evaluation of the limiting coverage.}
Normal symmetry and the critical-value identity give
\[
\Phi(-z_{1-\gamma/2})
=1-\Phi(z_{1-\gamma/2})
=\gamma/2.
\]
The standard normal law has zero mass at the interval endpoints. Hence
\[
N(0,1)([-z_{1-\gamma/2},z_{1-\gamma/2}])
=
\Phi(z_{1-\gamma/2})-\Phi(-z_{1-\gamma/2})
=
1-\gamma.
\]
Substitution into the Wald-event limit proves
\[
\Pr_\theta\left\{
|\widehat\tau^*-\tau|
\le z_{1-\gamma/2}\sqrt{\widehat V^*/n}
\right\}\longrightarrow1-\gamma,
\]
completing all the asserted conclusions.
% lean: CausalSmith.Stat.LdpAteEfficiencySurface.gaussianMeasure_waldCriticalValue_Icc
\end{enumerate}
\end{proof}
